%
% menu_presence.tex
%
% A Z specification of which sessions earn menu presence -- the agent
% ("Details") entries in the Display's Clients menu -- and of the two things
% every change to that set owes the rest of the Hub: a re-push of the menus
% to the Display, and exactly one bound departure sink for the session.
%
% WHY THIS MODEL EXISTS. The same defect class recurred in three review
% rounds of bead lux-ghhf: "a change to which sessions earn menu presence
% happens without the Display being re-pushed, or a re-push/departure hook is
% left dangling." Lux CLAUDE.md ("Formal Verification (z-spec)") makes a
% model-check mandatory at the second recurrence; this is it. The instances:
%
%   D1  identify recorded a new agent identity but never marked the menus.
%   D2  a listen-leg (ws_listen) agent attach never marked the menus.
%       (D1 and D2 are the same shape: a presence-adding entry point that
%       forgot to mark. One fidelity control, D1, stands for both.)
%   D3  a listen-leg session is lease-reaped while its socket stays open; the
%       next RenewFrame/SubscribeFrame calls HubClientRegistry.record(conn,
%       identity) and silently re-registers it: presence returns with no mark
%       and no departure sink re-bound, so the entry is missing and, when it
%       later departs, becomes a ghost the departure never marks away.
%   D4  an mcp-session agent's lease lapses before the 15 s sweep runs;
%       another client's push drops it (NamedSessions filters on liveness);
%       the agent's next call renews it; MenuArming.track compares identity
%       only (unchanged) so nothing marks and the entry stays missing.
%   D5  Operations.identify reads the "before" identity OUTSIDE the store
%       write lock; a departure landing between that read and
%       identify_client makes the re-registration look "unchanged" to track,
%       so no sink is bound and no mark is raised: a later ghost entry.
%   D6  ws_listen binds a fresh lambda as the departure sink on every attach;
%       DepartureSinks.bind dedups by equality and lambdas never compare
%       equal, so each reconnect adds one more sink, each pinning a dead
%       HubListenSession.
%
% All six are one violation of one rule, which the model states once: the
% presence set is a function of registry state, and EVERY registry
% transition that changes it must, in the same critical section, raise the
% menu mark and leave exactly one sink bound -- no caller is trusted to
% remember. The model is stated for the fixed design. The fidelity controls
% menu_presence_d{1,3,4,5,6}_buggy.tex each alter exactly one operation and
% make the corresponding invariant's negation reachable.
%
% THE DESIGN THE MODEL PROVES (see also the Design section below).
%   (a) Presence is the LEASE-FILTERED set (DES-058, DES-064 s.1, DES-057):
%       registered, an agent, and its lease live. A lapse is time passing
%       with no Hub event, so it cannot mark; the entry may linger until the
%       sweep (<= one sweep interval) or the next push. A renewal that
%       revives a lapsed session IS a presence transition and marks (D4);
%       the sweep departs every lapsed session and marks.
%   (b) The transitions that do change presence (absent to present, present
%       to absent) mark and bind/unbind under the Hub store's write lock, in
%       one step, and are all routed through HubDisplay (the outer lock);
%       the registry reports the change, it does not call out of its own
%       lock. menu_presence_locks.tex shows why: a registry that fired the
%       hook itself would invert the lock order and deadlock.
%   (c) The sink is ONE shared object per Hub, so a bind that repeats is a
%       no-op under DepartureSinks' equality dedup (D6).
%
% Read first: DESIGN.md DES-057 (identity, leases), DES-064 (the Clients
% menu), DES-067 (grouping), DES-094 (liveness is a heuristic), DES-098
% (menu dispatch / MenuDepartureSink), docs/connection_lease_reaping.tex
% (the departure cascade and its three paths; I1, I5, I7) and
% docs/menu_lifecycle.tex (MO: menuOwner within registered), and the code:
% src/punt_lux/domain/hub/hub_clients.py (record, attach_listener,
% reap_lapsed_locked, named_sessions), departure_sinks.py (bind, fire),
% hub_display.py (identify_client, _depart, _depart_lapsed, write_lock),
% operations/menu_arming.py (admit, track), operations/facade.py (identify),
% src/punt_lux/ws_listen.py (_attach_and_register_writer, _read_loop,
% _teardown).
%
% Type-check:  fuzz docs/menu_presence.tex
% Model-check: probcli docs/menu_presence.tex -model_check -p DEFAULT_SETSIZE 3
% The safety goals (P1, P2) are reachability checks of their negations; see
% the Verification section for the exact -goal predicates.
%

\documentclass[a4paper,10pt,fleqn]{article}
\usepackage[margin=1in]{geometry}
\usepackage{fuzz}
\usepackage[colorlinks=true,linkcolor=blue,citecolor=blue,urlcolor=blue]{hyperref}

\begin{document}

\title{lux Menu Presence: a Z Specification}
\author{Formal model of which sessions earn a Clients-menu entry, and of the
  mark and the departure sink every change to that set owes the Hub}
\date{October 2026}
\maketitle

% ============================================================
\section{Introduction}
% ============================================================

The Display's Clients menu lists the sessions that hold a lease and have
something to show: a session that registered a command, or an agent (an
\texttt{mcp-session}), which is kept present on its kind alone
(\texttt{ClientSession.earns\_menu\_presence}). The Hub does not push the menu
continuously. It raises a \emph{mark} (\texttt{DirtyMarker.mark\_menus}); the
replicator, on its next tick, reads the roster and sends the Display the
whole of it. A session's departure is announced the same way, by a
\emph{departure sink} bound in the Hub's \texttt{DepartureSinks} and fired
when the session leaves the registry.

Two obligations follow, and every defect in this class breaks one of them.

\begin{description}
  \item[P1 --- the Display converges.] When no mark is pending, no live
    session that earns presence is missing from the Display, and nothing the
    Display shows is unregistered; a lapsed entry may linger only until the
    sweep or the next push. A change to the presence set without a mark leaves
    the Display wrong until some unrelated event marks: a missing entry
    (D1--D5) or, worse, a ghost entry (D3, D5) that outlives its session.
  \item[P2 --- exactly one sink.] Every registered session that earns
    presence has exactly one departure sink bound, and no departed session
    has any. Zero sinks is the ghost: the departure fires nothing, so
    nothing marks the entry away. More than one is the leak (D6): each
    surplus sink pins a dead object for the life of the process.
\end{description}

Both are kept out of the state schema, as in
\texttt{connection\_lease\_reaping.tex}: a predicate in the state schema is
enforced as a guard on every successor, which silently disables the
operation that would violate it. The schema below carries structural
coherence only; P1 and P2 are checked as the unreachability of their
negations.

The model is flat and bounded: a carrier of connections, the registry's
membership and lease state, the count of bound sinks per connection, the
pending-mark flag, and the Display's shown set. Identity is abstracted to
the one fact that matters here: whether the session is an agent, hence earns
presence. Command-registering sessions earn presence by a second route
(\texttt{callbacks}); it is the same obligation on a different transition and
adds no behaviour this model does not already exhibit, so it is elided.

% ============================================================
\section{Basic Types and Free Types}
% ============================================================

\begin{zed}
  [CONN]
\end{zed}

A connection identity, the \texttt{ConnectionId} keying the registry, the
listen legs and the departure sinks. One identity is shared by successive
sessions of one agent (\texttt{connection\_for} derives it from the declared
identity), so a reconnect is the same \texttt{CONN} arriving again. Three
values suffice for an interleaving between two agents with a third bystander;
the carrier is bounded by ProB's \texttt{DEFAULT\_SETSIZE}.

\begin{zed}
  FSTATE ::= fclear | fset
\end{zed}

The pending-mark flag: no mark pending, or a mark raised and not yet
consumed by a push. Prefixed constructors, never \fuzz's logical
\texttt{true}/\texttt{false}.

% ============================================================
\section{State}
% ============================================================

\begin{schema}{MenuPres}
  registered : \power CONN \\
  earns : \power CONN \\
  lapsed : \power CONN \\
  open : \power CONN \\
  sinks : CONN \fun 0 \upto 2 \\
  dirty : FSTATE \\
  shown : \power CONN
\where
  earns \subseteq registered \\
  lapsed \subseteq registered
\end{schema}

\texttt{registered} is the registry's membership
(\texttt{HubClientRegistry.\_sessions}), live or lapsed. \texttt{earns} is
the registered sessions that earn presence: agents. \texttt{lapsed} is the
registered sessions whose lease has expired and which no sweep has yet
removed; it records the passage of time and nothing else. \texttt{open} is
the connections whose listen socket is still open, which the registry does
not know and which is why a session can be swept while its socket lingers
(D3). \texttt{sinks} counts the departure sinks bound per connection,
saturating at two: two is already the leak, and a larger count would only
enlarge the state space. \texttt{dirty} is the pending mark. \texttt{shown}
is the set the Display currently shows.

\begin{schema}{Init}
  MenuPres'
\where
  registered' = \emptyset \\
  earns' = \emptyset \\
  lapsed' = \emptyset \\
  open' = \emptyset \\
  sinks' = CONN \cross \{ 0 \} \\
  dirty' = fclear \\
  shown' = \emptyset
\end{schema}

% ============================================================
\section{Operations}
% ============================================================

\subsection{Recording an agent --- one transition, three entry points}

\texttt{AgentRecord} is the single registry transition by which a session
that is an agent is recorded: the upsert \texttt{HubClientRegistry.record}
performs for \texttt{identify\_client}, for \texttt{attach\_listener}, and for
the read loop's per-frame \texttt{record(conn, identity)}. It registers the
connection if absent, renews its lease, and adds it to the presence set. If
that changed the presence set (the connection was not already earning, or
was earning but lapsed, so the record revives it) it raises the mark; in
every case it leaves exactly one sink bound. Everything
happens in one step, because the fixed design performs it under the store's
write lock (Section~\ref{sec:design}); the registry reports whether presence
changed and the caller of record is not asked to compare anything.

\begin{schema}{AgentRecord}
  \Delta MenuPres \\
  c? : CONN
\where
  registered' = registered \cup \{ c? \} \\
  earns' = earns \cup \{ c? \} \\
  lapsed' = lapsed \setminus \{ c? \} \\
  sinks' = sinks \oplus \{ c? \mapsto 1 \} \\
  (c? \in earns \land c? \notin lapsed \land dirty' = dirty) \lor
    ((c? \notin earns \lor c? \in lapsed) \land dirty' = fset) \\
  shown' = shown
\end{schema}

\begin{zed}
  Identify \defs AgentRecord \land [\Delta MenuPres | open' = open]
\end{zed}

\texttt{Identify} models \texttt{Operations.identify}: the declaration is
recorded and, if it newly earns presence, marked and sunk, in one step.
This is the D1 and D5 fix: there is no \texttt{before} snapshot to go stale,
because the comparison that decides the mark is made against the registry's
own state inside the transition.

\begin{zed}
  Attach \defs AgentRecord \land [\Delta MenuPres; c? : CONN | open' = open \cup \{ c? \}]
\end{zed}

\texttt{Attach} models \texttt{HubListenSession.\_attach\_and\_register\_writer}:
the same transition plus the socket being open. This is the D2 and D6 fix: the
sink bound is the single shared one, so attaching again leaves \texttt{sinks}
unchanged (overriding the count with one is idempotent).

\begin{zed}
  ReadLoopRecord \defs AgentRecord \land
    [\Delta MenuPres; c? : CONN | open' = open \land c? \in open]
\end{zed}

\texttt{ReadLoopRecord} models the per-frame \texttt{record(conn, identity)}
in \texttt{HubListenSession.\_read\_loop}, reachable only while the socket is
open. When the session is registered it is a renewal and changes nothing but
\texttt{lapsed}. When the session has been swept it is a revival, the D3 path,
and is by construction the same transition as a first attach: mark and sink.

\subsection{Contact and time}

\begin{schema}{Connect}
  \Delta MenuPres \\
  c? : CONN
\where
  registered' = registered \cup \{ c? \} \\
  earns' = earns \\
  lapsed' = lapsed \setminus \{ c? \} \\
  open' = open \\
  (c? \in earns \land sinks' = sinks \oplus \{ c? \mapsto 1 \}) \lor
    (c? \notin earns \land sinks' = sinks) \\
  (c? \in earns \land c? \in lapsed \land dirty' = fset) \lor
    (\lnot (c? \in earns \land c? \in lapsed) \land dirty' = dirty) \\
  shown' = shown
\end{schema}

A contact that is not an identify (\texttt{register\_client}, a REST call):
registers and renews. For a session that earns nothing it owes nothing; for
an earning session whose lease had lapsed it is a revival, hence a presence
transition: it marks, and leaves the one sink bound.

\begin{schema}{Renew}
  \Delta MenuPres \\
  c? : CONN
\where
  c? \in registered \\
  registered' = registered \\
  earns' = earns \\
  lapsed' = lapsed \setminus \{ c? \} \\
  open' = open \\
  (c? \in earns \land sinks' = sinks \oplus \{ c? \mapsto 1 \}) \lor
    (c? \notin earns \land sinks' = sinks) \\
  (c? \in earns \land c? \in lapsed \land dirty' = fset) \lor
    (\lnot (c? \in earns \land c? \in lapsed) \land dirty' = dirty) \\
  shown' = shown
\end{schema}

Any contact from a registered session renews its lease
(\texttt{renew\_if\_registered}, the \texttt{\_renewed} step of \texttt{record}).
A renewal of a lapsed earning session puts it back in the lease-filtered
presence set: it marks and re-asserts the one sink (D4). A renewal of a live
session changes nothing the Display is told about.

\begin{schema}{LeaseLapse}
  \Delta MenuPres \\
  c? : CONN
\where
  c? \in registered \\
  c? \notin lapsed \\
  registered' = registered \\
  earns' = earns \\
  lapsed' = lapsed \cup \{ c? \} \\
  open' = open \\
  sinks' = sinks \\
  dirty' = dirty \\
  shown' = shown
\end{schema}

Time passes with nothing to renew the lease. It leaves the presence set
(a lapsed session is no longer live) but it is not a Hub event, so nothing can
mark: the Display may keep showing the entry until the sweep departs it or the
next push re-reads the roster. P1 states exactly how long this may last
(Section~\ref{sec:inv}).

\subsection{Departure}

\begin{schema}{Departure}
  \Delta MenuPres \\
  c? : CONN
\where
  c? \in registered \\
  registered' = registered \setminus \{ c? \} \\
  earns' = earns \setminus \{ c? \} \\
  lapsed' = lapsed \setminus \{ c? \} \\
  open' = open \\
  sinks' = sinks \oplus \{ c? \mapsto 0 \} \\
  (c? \in earns \land dirty' = fset) \lor
    (c? \notin earns \land dirty' = dirty) \\
  shown' = shown
\end{schema}

The one deregister-and-release step (\texttt{HubDisplay.\_depart}): the
session leaves the registry, its sinks fire and are unbound, and the sink
marks the menus. Note \texttt{open'}$=$\texttt{open}: the registry has no
authority over the socket.

\begin{zed}
  DropGraceful \defs Departure
\end{zed}

\begin{zed}
  Sweep \defs Departure \land [MenuPres; c? : CONN | c? \in lapsed]
\end{zed}

\texttt{DropGraceful} is \texttt{drop\_connection}: no lease condition.
\texttt{Sweep} is the \texttt{LeaseReapSweep} / \texttt{\_depart\_lapsed}
path, guarded on a lapsed lease. Both funnel through the same transition,
which is \texttt{connection\_lease\_reaping.tex}'s I5.

\begin{schema}{Teardown}
  \Delta MenuPres \\
  c? : CONN
\where
  c? \in open \\
  registered' = registered \\
  earns' = earns \\
  lapsed' = lapsed \\
  open' = open \setminus \{ c? \} \\
  sinks' = sinks \\
  dirty' = dirty \\
  shown' = shown
\end{schema}

The socket closes (\texttt{HubListenSession.\_teardown}). An agent's presence
does not depend on its leg, so the session, its identity, its lease and its
sink persist; a reconnect within the lease finds them.

\subsection{The replicator push}

\begin{schema}{Push}
  \Delta MenuPres
\where
  dirty = fset \\
  registered' = registered \\
  earns' = earns \\
  lapsed' = lapsed \\
  open' = open \\
  sinks' = sinks \\
  dirty' = fclear \\
  shown' = earns \setminus lapsed
\end{schema}

The replicator consumes a pending mark: it clears the flag and sends the
Display the whole lease-filtered presence set
(\texttt{NamedSessions.over(\_live\_locked(), ...)}). Atomic here for the reason a standard
dirty-flag protocol is: the flag is cleared before the roster is read, so a
mark raised after the read survives for the next tick.

% ============================================================
\section{Invariants}
\label{sec:inv}
% ============================================================

\paragraph{P1 --- the Display converges.}
\[
  dirty = fclear \;\Rightarrow\;
    (earns \setminus lapsed) \subseteq shown \;\land\; shown \subseteq earns .
\]
Presence is the lease-filtered set $earns \setminus lapsed$ (DES-058, DES-064
\S1, DES-057). With no mark pending, (i) no live earning session is ever
missing from the Display, and (ii) the Display shows nothing that is not a
registered earning session, so no ghost. What it may show beyond the live set is
exactly a \emph{lapsed, unswept} entry that was shown when the lease ran out:
$shown = (earns \setminus lapsed) \cup (shown \cap lapsed)$. A lapse is time
passing with no Hub event, so nothing can mark it; the entry lingers until the
sweep departs it (which marks) or the next push re-reads the roster, never
longer. The bound is the \emph{sweep obligation}: \texttt{Sweep(c)} is enabled
for every lapsed registered $c$ and departs it and marks, so the lingering set
$lapsed$ is emptied by operations the Hub's timer is bound to run
(\texttt{connection\_lease\_reaping.tex} I1). Every transition that adds to
$earns \setminus lapsed$ (\texttt{AgentRecord}, and \texttt{Connect} or
\texttt{Renew} on a lapsed earning session) sets \texttt{dirty} in the same
step; every transition that removes from \texttt{earns} (\texttt{Departure})
does too; \texttt{Push} sets \texttt{shown} to the live set while clearing
the flag. The negation, checked as a goal, is
$dirty = fclear \land (\lnot((earns \setminus lapsed) \subseteq shown) \lor
\lnot(shown \subseteq earns))$.

\paragraph{P2 --- exactly one sink.}
\[
  (\forall c : CONN @ c \in earns \Rightarrow sinks(c) = 1)
  \;\land\; (\forall c : CONN @ sinks(c) \leq 1)
  \;\land\; (\forall c : CONN @ c \notin registered \Rightarrow sinks(c) = 0) .
\]
An earning session has one sink; no session has more than one; a departed
session has none. The first conjunct is the ghost (D3, D5), the second the leak
(D6), the third the dangling hook.

\paragraph{Why P1 has this shape.}
The roster the Display shows is read by
\texttt{NamedSessions.over(\_live\_locked(), ...)}, which filters out sessions
whose lease has lapsed, as the ADRs require. A lease passing is therefore a
presence change that no Hub transition performs. P1 cannot demand equality
with the live set at every quiescent instant; it demands the strongest thing
that is true: nothing live is missing, nothing unregistered is shown, and a
lapsed entry's lingering is cut off by the sweep. D4 is the renewal that ends
a lapse without marking: the live set grows, the flag stays clear, and conjunct
(i) fails \emph{for good}, since no later event is owed to repair it.

% ============================================================
\section{The Design the Model Proves}
\label{sec:design}
% ============================================================

\paragraph{Presence is lease-filtered; a revival is a transition.}
A session earns presence while it is registered, an agent, and its lease is
live: DES-058 (``entries live on the session's lease and leave when it
lapses''), DES-064 \S1 (``an entry exists exactly while its registrant holds a
lease''), DES-057 (``leases expire state whose owner left''). No ADR is
amended. The consequences the model makes explicit: a lapse is time, not an
event, so it cannot mark and the entry may linger until the sweep or the next
push; a renewal that ends a lapse \emph{is} an event that grows the presence
set, so it must mark and re-assert the one sink (\texttt{Renew}, \texttt{Connect},
and \texttt{AgentRecord} on a lapsed session); and the sweep must depart every
lapsed session and mark. D4 is the middle one forgotten.

\paragraph{The registry owns the transition; the caller owns nothing.}
\texttt{AgentRecord} is one step, so a caller cannot forget the mark or bind
a sink twice. In code the step is a \texttt{HubDisplay} method (the outer
\texttt{StoreLock}) that calls a registry upsert returning whether presence
changed, then marks and binds the one shared sink. All three entry points
call it: \texttt{identify}, \texttt{attach}, and the read loop's
\texttt{record}. The read loop must no longer call
\texttt{HubClientRegistry.record} directly.

\paragraph{The registry does not call out of its own lock.}
The tempting shape is for \texttt{HubClientRegistry} itself to raise the mark
and bind the sink inside \texttt{record}. It cannot: binding a sink takes the
store's write lock, and the registry lock is taken \emph{inside} the store
lock on every existing path (\texttt{HubDisplay.register\_client},
\texttt{\_depart}). A registry that took the store lock while holding its own
would invert that order. \texttt{menu\_presence\_locks.tex} proves the
inversion deadlocks (two threads, one in each order) and that the shape above
does not.

\paragraph{One shared sink.}
The sink is a single object constructed once (the production
\texttt{MenuDepartureSink} already is; \texttt{ws\_listen} binds a per-attach
lambda). With one object, \texttt{DepartureSinks.bind}'s equality dedup makes
the repeated bind a no-op, which is the \texttt{oplus} to one above, and
\texttt{\_teardown} need not unbind it: the sink belongs to the session, not
to the socket.

% ============================================================
\section{Verification}
\label{sec:verify}
% ============================================================

Type-checking is a \texttt{fuzz} pass:
\begin{verbatim}
  fuzz docs/menu_presence.tex
\end{verbatim}

P1 and P2 are checked by asking ProB whether their negation is reachable. A
goal reported \emph{not found} means the invariant holds on every reachable
state.

\begin{verbatim}
  # P1 (goal NOT found in this, the fixed, spec)
  probcli docs/menu_presence.tex -model_check -nodead \
    -goal "dirty = fclear & (not((earns \ lapsed) <: shown) or not(shown <: earns))" \
    -p DEFAULT_SETSIZE 3

  # P2 (goal NOT found in this, the fixed, spec)
  probcli docs/menu_presence.tex -model_check -nodead \
    -goal "#c.(c : CONN & ((c : earns & sinks(c) /= 1) or sinks(c) > 1 \
      or (c /: registered & sinks(c) /= 0)))" -p DEFAULT_SETSIZE 3
\end{verbatim}

The positive outcomes are confirmed reachable, so the invariants are not
vacuously true of a model that never shows an entry:

\begin{verbatim}
  # an entry is shown with no mark pending (FOUND)
  probcli docs/menu_presence.tex -model_check -nodead \
    -goal "dirty = fclear & shown /= {}" -p DEFAULT_SETSIZE 3

  # the permitted lingering: a lapsed, unswept agent is still shown with no
  # mark pending (FOUND -- the honest cost of a lapse being time, not an event)
  probcli docs/menu_presence.tex -model_check -nodead \
    -goal "dirty = fclear & #c.(c : shown & c : lapsed)" \
    -p DEFAULT_SETSIZE 3

  # a live agent is shown, nothing lapsed, no mark pending (FOUND)
  probcli docs/menu_presence.tex -model_check -nodead \
    -goal "dirty = fclear & lapsed = {} & shown /= {} & earns = shown" \
    -p DEFAULT_SETSIZE 3
\end{verbatim}

Deadlock-freedom and the structural invariant over the full reachable state
space:
\begin{verbatim}
  probcli docs/menu_presence.tex -model_check -p DEFAULT_SETSIZE 3
\end{verbatim}

The fidelity controls, each differing from this specification in one
operation:

\begin{verbatim}
  menu_presence_d1_buggy.tex   P1 negation FOUND (identify never marks)
  menu_presence_d3_buggy.tex   P1 and P2 negations FOUND (revival by record)
  menu_presence_d4_buggy.tex   P1 negation FOUND (renewal of a lapsed
                               agent does not mark)
  menu_presence_d5_buggy.tex   P1 and P2 negations FOUND (stale before-read)
  menu_presence_d6_buggy.tex   P2 negation FOUND (a sink per attach)
\end{verbatim}

The lock order is checked in the companion models
\texttt{menu\_presence\_locks.tex} (deadlock-free) and
\texttt{menu\_presence\_locks\_inverted\_buggy.tex} (deadlock found).

Re-run \texttt{fuzz} and these goals whenever
\texttt{hub\_clients.py}'s \texttt{record}/\texttt{attach\_listener}/
\texttt{reap\_lapsed\_locked}/\texttt{named\_sessions},
\texttt{departure\_sinks.py}, \texttt{hub\_display.py}'s
\texttt{identify\_client}/\texttt{\_depart}/\texttt{\_depart\_lapsed},
\texttt{operations/menu\_arming.py}, \texttt{operations/facade.py}'s
\texttt{identify}, or \texttt{ws\_listen.py} change.

\end{document}
